\section{Auxiliary Clustering Experiment}\label{app:aux-rep}

This appendix records a representation-learning experiment that we treat as descriptive rather than as core evidence \citep{tishby1999bottleneck,geiger2015optimalkullbackleibleraggregation}. It asks what happens when the packaging is learned from noisy observations instead of being given.

\paragraph{Design.} We simulate $3{,}000$ steps of the hidden-type chain used in the budget benchmark (\Cref{sec:results-budget}), after a burn-in of $200$ steps. Each microstate emits a four-dimensional Gaussian observation whose mean depends mostly on its true package, with a small microstate-specific offset. We fit $k$-means \citep{macqueen1967somemethods} with $k\in\{2,4,8\}$ to the observations in the first half of the sequence and label every step with its nearest cluster. That label sequence is the learned packaging $Y_t$. First- and second-order count predictors are fitted to the labels in the first half and scored on the second half, using the same $1{,}498$ targets for both orders. Finally, we compute the plug-in estimate
\[
\widehat I(X_t;Y_{t+1}\mid Y_t)=\widehat H(Y_{t+1}\mid Y_t)-\widehat H(Y_{t+1}\mid X_t,Y_t)
\]
from transition counts over the whole sequence, where $X_t$ is the true hidden microstate.

Two features distinguish this setting from the main text. First, the labels come from noisy observations, so $Y_t$ is not a function of $X_t$. The second entropy above therefore conditions on $(X_t,Y_t)$, and the decomposition of \Cref{prop:cd-decomposition} does not apply. Second, the estimate is a raw plug-in statistic. Its upward bias grows with the number of contexts and outcomes, here $8k$ contexts and $k$ outcomes estimated from $2{,}999$ transitions.

\paragraph{Results.} As $k$ increases from $2$ to $4$ to $8$, the plug-in estimate rises from $0.078$ to $0.136$ to $0.158$ nats (\Cref{fig:appendix-rep}a). Held-out log loss rises too (\Cref{fig:appendix-rep}b): from $0.626$ to $1.307$ to $2.005$ nats for order one, and from $0.628$ to $1.318$ to $2.837$ for order two.

\begin{figure}[tbp]
\centering
\includegraphics[width=\linewidth]{fig_clustering}
\caption{Auxiliary clustering experiment. (a)~Plug-in estimate of $\widehat I(X_t;Y_{t+1}\mid Y_t)$ for learned $k$-means packagings. (b)~Held-out log loss of first- and second-order predictors on the learned labels, with the uniform-guess loss $\log k$ for scale. Because the predicted alphabet grows with $k$, losses at different $k$ are not directly comparable.}
\label{fig:appendix-rep}
\end{figure}

\paragraph{Interpretation.} These trends should be read with care. The rising loss mostly reflects the growing alphabet: predicting one of eight labels is harder than predicting one of two, and the order-one loss stays $0.07$ to $0.08$ nats below the uniform-guess value $\log k$ at every $k$. The order-two loss at $k=8$ exceeds $\log 8$, a sign of overfitting $64$ contexts with $1{,}500$ training steps. The rising plug-in estimate is consistent with learned clusters that split the data in ways the dynamics do not respect. Plug-in bias, which also grows with $k$, could contribute to it as well, and this experiment does not separate the two. The conclusion we draw is therefore modest. More clusters did not reduce the estimated deficit here. That more refinement need not help is established exactly by \Cref{ex:four-cycle}, not by this experiment.

We also ran exploratory neural bottleneck encoders. In those runs, the learned codes used fewer categories than their nominal size, so they do not give a clean comparison at equal budget. We do not use them as evidence.
